% 0.摘要.tex —— 摘要
\begin{center}
  {\LARGE\bfseries 药材烘干过程温度与水分浓度分布的建模与数值求解}\\[6pt]
  {\large ——基于守恒离散与变物性迭代的一维热质传递模型}
\end{center}

\begin{abstract}
\noindent 针对半径 $2$ cm、长 $25$ cm 的圆柱形药材，本文采用径向一维非稳态模型，依次研究预热期温度与含水率分布、全烘干过程演化、全场含水率达标时间，以及考虑尺寸收缩的烘干时长。该模型描述长圆柱中段的径向传递，忽略端部效应和汽化潜热（前三个问题还忽略收缩）；相关结果均以这些简化和题设经验参数为前提。

对于问题一，采用附录 2 的常热物性与 $D(C)$ 扩散关系，热、质方程相互解耦。对中心、内部和表面节点分别作守恒离散，以显式 Euler 为基线，采用 Crank--Nicolson（CN）与 Thomas 追赶法求解。$1800$ s 时中心与表面温度分别为 $33.5752$、$36.7853$ $^\circ$C，表面含水率为 $1.5120$ kg/kg；表面与体积平均含水率分别下降约 $40.70\%$、$10.11\%$，表明短时失水集中于外层。$\Delta t=1$ s 时，两种格式终态最大差为 $5.02\times10^{-3}$ $^\circ$C 与 $8.74\times10^{-5}$ kg/kg，故四位小数输出不代表两种方法具有相同数值精度。

对于问题二，采用含水率相关的 $\rho,c_p,k$ 与温度、含水率相关的 $D(C,T)$，形成经物性耦合的非线性系统。热方程采用局部等效热容近似；水分方程含扩散系数敏度项的 Newton 型线性化与热场块交替更新相结合。环境曲线在前 $4$ h 内插值，之后保持末值，计算覆盖 $72$ h。$3$ h 时径向温差约为 $0.12$ $^\circ$C，而中心与表面含水率仍为 $1.7662$、$1.0078$ kg/kg，显示热响应与失水过程具有不同时间尺度。

对于问题三，沿用变物性模型，以 $\max_r C(r,t)<0.15$ kg/kg 为全场判据，结合时间步加倍估计、稠密输出与重新积分二分定位求达标时间。终止时间按三种口径给出并明确区分：交付网格 $N=640$ 的\textbf{数值值}为 $57.1588$ h；综合均匀与分级网格的\textbf{外推趋势估计}为 $57.17$ h（$57.1696\sim57.1723$ h）；按整数分钟向上取整得到的\textbf{保守生产时间}为 $3431$ min $=57.1833$ h（$57$ 小时 $11$ 分）。三者中，$57.1588$ h 是单一网格的直接输出，$57.17$ h 是网格外推的趋势估计（非严格误差界），$57.1833$ h 是面向生产调度的保守取整。网格差异是所考察数值设置下终止时间的主要误差来源；该结果不构成实际干燥时间的秒级精度保证。

对于问题四，考虑烘干过程中的尺寸收缩，引入附件 2 的收缩半径经验数据 $R(t)$（$2.0\to1.198$ cm）作为移动边界，采用 Landau 坐标变换将移动域映射为固定域，在附录 4 变物性下用含对流项的守恒型 Crank--Nicolson 格式求解。严格链式法则给出的对流项符号为正（题面符号为笔误），符号敏感性检验确认该对流项使烘干时长增大 $1.57$ h、约占 $3\%$。烘干时长为 $188640$ s $=52.40$ h（$3144$ min，整数分钟向上取整），较冻结半径 $2$ cm 的 $129.1$ h 缩短 $2.46$ 倍，表明收缩是烘干时长的决定性因素；空间（$N=400/800$）与时间（$\Delta t=60/120$ s）收敛均达秒级一致。

\noindent\textbf{关键词：}径向热质传递；变物性模型；守恒离散；Newton 型线性化；自适应时间步；网格收敛；移动边界；Landau 变换
\end{abstract}
